\begin{papercard}{dehling2015}
\cf{Stages} \stagetag{S} \stagetag{A} \stagetag{H}. A store-metadata-only overview of about 24{,}000 English mHealth apps. Apps were clustered by tagged functionality and rated on five \emph{potential damage} characteristics, yielding 12 app archetypes. This is the main precedent for a \emph{functional} typology linked to harm potential.

\cf{Aim} An overview of mHealth apps on iOS and Android focused on \emph{potential damage to users} through security and privacy infringements, treating apps as diverse rather than monolithic \pc{dehling2015}{2}{research is to establish an overview of mHealth apps offered}.

\cf{Discovery} Medical and Health \& Fitness categories of both stores in May 2013. iOS was parsed from HTML. Android was searched through a third-party API, but Google returns at most 500 apps per query, so queries were built from a 111{,}632-entry English word list, each word appended with ``medical'' and with ``health'' \pc{dehling2015}{2}{only a maximum of 500 apps is returned per search request}.

\cf{Tagging filter} Tags were attached to \emph{strings} in descriptions (e.g.\ medication, pharmaceutical and drug map to one tag), building a reusable string-to-tag corpus for automatic tagging \pc{dehling2015}{3}{Instead of assigning tags directly to an app, we assigned tags to}. Cut-offs were agreed in author discussions because no guidance existed. 200 frequently rated apps were tagged manually; the rest were tagged by string matching, and apps with fewer than four distinct tags were excluded \pc{dehling2015}{3}{cut-off points were determined according to the available data}. 37{,}246 discovered, 12{,}841 excluded \pc{dehling2015}{4}{We discovered a total of 37,246 apps}.

\cf{Clustering} An app--tag graph (vertices are apps, edges are shared tags) clustered with the Louvain modularity heuristic, which needs no preset number of clusters \pc{dehling2015}{3}{For identification of clusters, we used a heuristic by Blondel}. Result: 245 clusters (modularity 0.47); 70 uninformative clusters were excluded \pc{dehling2015}{7}{Application of the Louvain method}.

\cf{Damage assessment scheme} Five characteristics from the health-IT security literature (Table 1): (1) \emph{specificity} of available information (standard, non-standard, medical); (2) damage through \emph{leaks}; (3) through \emph{manipulation}; (4) through \emph{loss}; (5) \emph{value to third parties}, each none/low/high \pc{dehling2015}{3}{we assess information security and}. Leak consequences are spelled out: socioeconomic repercussions, embarrassment and reputation, stigma, loss of family respect, medical fraud and identity theft, insurance and employment effects \pc{dehling2015}{3}{such information increase the likelihood of potential damage}. Third-party value: insurance and employee selection, targeted marketing, celebrity data \pc{dehling2015}{4}{Further uses of others}. Standard data (location, device ID) deliberately did not raise the rating, which focuses on mHealth-specific risk \pc{dehling2015}{3}{Cluster assessment is focused on risks specific}.

\cf{Rating procedure and reliability} Two researchers assessed every cluster, first writing a short description from tag prevalence and verifying it against member apps. Janson--Olsson $\iota=0.71$ (substantial). Differences were discussed, with a third researcher if needed \pc{dehling2015}{4}{There were two researchers that assessed all discovered clusters}. Clusters with identical ratings were merged into archetypes \pc{dehling2015}{4}{are identified by grouping clusters with}.

\cf{Results} Most apps are rarely rated: 81\,\% of iOS and 76\,\% of Android apps had fewer than 10 ratings, so less than a quarter are in widespread use \pc{dehling2015}{5}{apps that have been rated less than 10 times}. Rating count proxies downloads on Android (Spearman 0.89) \pc{dehling2015}{5}{For Android apps, rating count and download count are strongly}. Twelve archetypes (Table 3), e.g.\ AT 5 \emph{Fitness Trackers} (non-standard data, low on all), AT 9 \emph{Intimate Trackers} (menstruation, intercourse, fertility, pregnancy: high leak damage), AT 12 \emph{Health Records} (high on all) \pc{dehling2015}{10}{Menstruation, intercourse, fertility, and pregnancy tracker}. 95.6\,\% pose some potential damage. Even \emph{installation} of some apps reveals sensitive information \pc{dehling2015}{12}{or the sole fact that a guide for stress relief}.

\cf{Limitations stated} Store information only, no installation or testing; Android coverage unknown because of the 500-result cap \pc{dehling2015}{12}{it was unfeasible to install and test all apps}.

\cf{Use for this thesis} For \stagetag{S}: a ready functional typology that distinguishes \emph{trackers} (fitness and intimate) from information providers and calculators, useful for defining ``health tracker'' by function rather than store category. For \stagetag{H}: the five-dimension damage scheme (leak, manipulation, loss, third-party value, specificity) and the explicit consequence list are a citable seed for consequence nodes, and the ``installation reveals'' point is a capability without data collection. For \stagetag{A}: query construction to bypass result caps, and ratings as a usage proxy. Out of date for store mechanics (2013).
\end{papercard}
